\scanpage{176}
% archive sheet number: 175
\deleted{de plus}, \textsuperscript{(1)} une somme de \deleted{tels} coefficients de
représentations, dans $V$, resp $W$ est un
coefficient d'une représentation dans $V\oplus W$. ---

\deleted{De plus, les}
Si on était \uline{parti} de $A_0$, puis passé à $\unsure{\mathscr G|\mathcal G|G}$,
on \unsure{retrouve son|retombe sur} $A_0$. Prouvons que partant de
$\mathscr G$, construisant $A_0$, alors \deleted{toute} représentation
de dim finie de $\Omega$ dont les coefficients
sont dans $A_0$, est dans $\mathscr G$. \deleted{\unsure{Démonstration|Dire}}
\deleted{par récurrence sur la dimension de $V$ :}
\deleted{Si $\dim V=1$, alors}

\uline{Lemme} 1 Toute représentation $\in\mathscr G(A_0)$ est
contenue dans une somme directe finie
représentations contenues dans $A_0$ (où
$\Omega$ opère par \unsure{cogénération|coopération|opération}) (On prend
une base $(e'_i)$ de $V'$, et on
prend la somme des représentations $x\unsure{-|\mapsto|\to} f_{x,e'_i}$)

Pour prouver $\deleted{\unsure{\mathscr G\subset}}\mathscr G(A_0)\subset\mathscr G$ on est ramené
au cas d'une représentation \uline{contenue}
dans $A_0$, à prouver qu'elle est $\in\mathscr G$.
Pour ceci, on note que \uline{lemme 2} toute
représentation \inserted{de dim finie} contenue dans $A_0$ est quotient
d'une somme \inserted{finie} de repr. \uline{monogènes}
contenues dans $A_0$ (trivial).

On est donc ramené au cas d'une
représentation dans $\Omega*f$, où $f\in A_0$.
On a $f=f_{xx'}$, $x,x'\in V$, $V$ \unsure{espace} \deleted{\unsure{de|d'une}}
